\appendix

\section{反常}\label{anomaly}

为说明从经典力学过渡到量子力学时的反常，我们考虑 $\R^2$，用泊松括号 $\{x_1,x_2\}=1$ 的实变量 $x_1,x_2$ 参数化。量子化后，$x_1$ 与 $x_2$ 成为满足 $[\h x_1,\h x_2]=-\i\hbar$ 的算符 $\h x_1,\h x_2$。在分析量子化问题时，我们假设 $x_i$ 的任一 $k$ 次多项式（$k\ge1$）的量子化应是 $\h x_i$ 的 $k$ 次多项式，且首项相同。问题在于：能否选择量子算符中因子的次序以及可能的低阶项，使换位子与泊松括号精确相符。

单项式 $x_{i_1}x_{i_2}\cdots x_{i_k}$ 的 Weyl 排序 $\left\langle\h x_{i_1}\h x_{i_2}\cdots\h x_{i_k}\right\rangle$ 定义为对因子 $k!$ 种可能次序的平均：
\be\label{weyl}\left\langle\h x_{i_1}\h x_{i_2}\cdots\h x_{i_k}\right\rangle=\frac{1}{k!}\left(\h x_{i_1}\h x_{i_2}\cdots\h x_{i_k}+\h x_{i_2}\h x_{i_1}\cdots\h x_{i_k}+\cdots\right).\ee
我们一开始并不假设 Weyl 排序是给因子排序的最佳方式。我们将看到，无论怎样排列因子、添加怎样的低阶项，都无法避免 $x_1,x_2$ 的多项式量子化中的反常。在下面将解释的意义上，Weyl 排序最接近避免反常。

对次数 $\le2$ 的多项式没有反常。次数 $\le2$ 的 Weyl 排序多项式的换位子，按预期方式与相应经典多项式的泊松括号相符。Weyl 排序是具有这一性质的唯一做法：给次数 $\le2$ 的 Weyl 排序多项式添加低阶项会改变其换位子。我们假设这一点已知（或读者会自行验证），并说明若考虑次数大于 $2$ 的多项式，就会出现反常。

$x_1$ 与 $x_2$ 的齐二次多项式，即 $x_1^2$、$x_2^2$ 与 $x_1x_2$，通过泊松括号生成 Lie 代数 $\SL(2,\R)$。$x_1,x_2$ 的 $k$ 次齐次多项式张成 $k+1$ 维向量空间，并按 $\SL(2,\R)$ 的 $k+1$ 维不可约表示变换。由于二次多项式无反常，$\h x_1,\h x_2$ 的 Weyl 排序二次多项式，即 $\h x_1^2$、$\h x_2^2$ 与 $\frac12(\h x_1\h x_2+\h x_2\h x_1)$，也通过换位子生成 $\SL(2,\R)$ 的 Lie 代数。$\h x_1,\h x_2$ 的 $k$ 次 Weyl 排序齐次多项式，像其经典对应物一样，按 $\SL(2,\R)$ 的 $k+1$ 维不可约表示变换。任何其他次序选择（等价地，给 $k$ 次 Weyl 排序多项式添加任何低阶项）都会破坏这一事实，因为次数小于 $k$ 的多项式按 $\SL(2,\R)$ 的较小表示变换。因此，当且仅当使用 Weyl 排序时，齐二次多项式与任意 $k$ 次齐次多项式的换位子才无反常。

然而，两个次数大于 $2$ 的 Weyl 排序多项式的换位子会出现反常。为说明这一点，只需在一个特殊情形中展示反常。令 $z=(x_1-\i x_2)/\sqrt2$、$\bar z=(x_1+\i x_2)/\sqrt2$，则 $\{z,\bar z\}=\i$ 且 $[\h z,\h{\bar z}]=\hbar$。我们有 $\{z^3,\bar z^3\}=9\i z^2\bar z^2$，故若无反常，就应要求
$[z^3,\bar z^3]=9\hbar\left\langle\h z,\h z,\h{\bar z},\h{\bar z}\right\rangle$。但简单计算表明，实际有
\be\label{tofo}[\h z^3,\h{\bar z}^3]=9\hbar\left\langle\h z,\h z,\h{\bar z},\h{\bar z}\right\rangle+\frac32\hbar^3,\ee
\begin{kanwu}[一处指数笔误]
原文式 \eqref{tofo} 左端写作 $[\h z^3,\h{\bar z}^2]$，与上下文（$\bar z^3$）不符；此处按文意改为 $\bar z^3$。
\end{kanwu}

这就展示了反常。

在这个特殊例子中，换位子里的反常项是 $c$-数。对更高次的多项式，比如 $z^3$ 与 $\bar z^4$，情况并非如此。这一断言无需进一步计算，由 Jacobi 恒等式
$[\h z,[\h z^3,\h{\bar z}^4]]+[\h{\bar z}^4,[\h z,\h z^3]]+[\h z^3,[\h{\bar z}^4,\h z]]=0$
即可推出。

\section{关于 $\Hom(\B_\cc,\B_\cc)$ 与 $\Hom(\B,\B_\cc)$ 的细节}\label{details}

\subsection{最低阶计算}

本附录的目的是概述 $\Hom(\B_\cc,\B_\cc)$ 与 $\Hom(\B,\B_\cc)$ 的计算，其中 $\B$ 是拉格朗日 $A$-膜。计算在 $\sigma$-模型微扰论的最低阶进行。可能的修正将在附录 \ref{corrections} 中讨论。

为具体定义复辛流形 $Y$ 的 $A$-模型（辛结构为 $\omega=\omega_K$），需要 $Y$ 上一个近复结构，使 $\omega_K$ 关于它为 $(1,1)$ 型且正。若 $Y$ 容许一个延拓其辛结构的超 Kähler 结构（意义见第 \ref{amo} 节），我们可以使用作为辛结构一部分的复结构 $K$。即使 $Y$ 不容许这样的超 Kähler 结构，也总能在 $Y$ 上取一个近超 Kähler 结构，即选取近复结构 $J$ 与 $K$，使它们（连同 $Y$ 的复辛结构）满足第 \ref{amo} 节讨论的代数关系。这之所以可能，理由如下。设 $Y$ 的实维数为 $4k$。$Y$ 的复辛结构把其切丛的结构群从 $\GL(4k,\R)$ 约化为 $\Sp(2k,\C)$。$\Sp(2k,\C)$ 的最大紧子群是紧形式 $\Sp_c(2k)$。商 $\Sp(2k,\C)/\Sp_c(2k)$ 可缩，故一旦切丛的结构群约化为 $\Sp(2k,\C)$，进一步约化到 $\Sp_c(2k)$ 就没有障碍。这正是定义 $J$、$K$ 以及 $Y$ 上满足通常代数关系的度量 $g$ 所需的结构。此外，任何两个这样的约化都同伦，因而给出等价的 $A$-模型。在下面的分析中，我们假设 $K$ 是某个近超 Kähler 结构的一部分。

首先考虑 $\Hom(\B_\cc,\B_\cc)$，沿用 \cite{GW} 第 2.2 节。考虑条带 $\Sigma=\R\times[0,\pi]$ 上的 $(\B_\cc,\B_\cc)$ 弦，第二个因子由 $\sigma$ 参数化。一般地，$\sigma$-模型的自由度是描述映射 $\Phi:\Sigma\to Y$ 的玻色场 $X$，以及取值于 $\Phi^*(TY)\otimes\SS_\pm$ 的费米场 $\psi_\pm$（$\SS_+$ 与 $\SS_-$ 是 $\Sigma$ 的正、负手征旋量丛）。$A$-模型微分 $Q$ 生成的 $X$ 的一般变换为
\be\label{windo}\delta X=(1-\i K)\psi_++(1+\i K)\psi_-.\ee
这意味着 $\delta X$ 有正比于 $\psi_+$ 的 $(1,0)$ 部分，以及正比于 $\psi_-$ 的 $(0,1)$ 部分。由于在 $\Sigma$ 内部，$\psi_+$ 与 $\psi_-$ 是 $\Phi^*(TY)\otimes\SS_\pm$ 的截面且不受特殊约束，为使 $\Sigma$ 内部的算符 $F(X,\bar X)$ 为 $Q$-不变，函数 $F$ 必须是常数。但我们关心的是插在边界上的局部算符，这里情况不同。一般地 \cite{ACNY}，在 CP 丛曲率为 $\sF$ 的秩 1 膜上、且存在二形式场 $\sB$ 时，世界面费米子满足边界条件
\be\label{wiffo}\psi_+=(g-(\sF+\sB))^{-1}(g+\sF+\sB)\psi_-.\ee
对 $\B_\cc$，$\sF+\sB=\omega_J$，边界条件化为
\be\label{nelm}\psi_+=J\psi_-,\ee
于是沿 $\Sigma$ 的边界，$X$ 的变分化为 $\delta X=(1+\i I)(1+J)\psi_-$。这意味着相对于复结构 $I$，沿 $\Sigma$ 的边界 $\delta X$ 只有 $(0,1)$ 部分：
\begin{align}\label{btrans}\delta^{1,0}X&=0,\cr
\delta^{0,1}X&=\rho,\end{align}
其中 $\rho=(1+\i I)(1+J)\psi_-$。由于 $Q^2=0$，式 \eqref{btrans} 表明沿 $\partial\Sigma$ 有
\be\label{bcorr}\delta\rho=0.\ee
在 $\rho$ 中次数为 $q$ 的边界算符 $F(X,\bar X,\rho)$ 对应于 $Y$ 上的 $(0,q)$-形式，而式 \eqref{btrans}、\eqref{bcorr} 表明 $Q$ 作用在这些算符上就如复结构 $I$ 中通常的 $\bar\partial$ 算符。因此在这个近似下，物理态空间是上同调 $\oplus_{q=0}^{\dim_\C Y}H^q(Y,\O)$。

现在考虑 $\Hom(\B,\B_\cc)$，其中 $\B$ 是秩 1 拉格朗日 $A$-膜，支撑在某个 $L\subset Y$ 上。为求 $Q$ 作用在 $(\B,\B_\cc)$ 弦上的上同调，先找出沿弦的玻色与费米零模。我们仍在同一批条带上工作，$\sigma=0$ 与 $\sigma=\pi$ 处的边界条件分别由 $\B_\cc$ 与 $\B$ 定义。在 $\sigma=\pi$ 处，映射 $\Phi:\Sigma\to Y$ 被约束为取值于 $L$。这一约束去掉了沿弦的一些玻色零模；剩下的零模描述映射 $\Phi:\partial\Sigma\to L$。对费米子，$\sigma=0$ 处的边界条件仍是 $\psi_+=J\psi_-$。在 $\sigma=\pi$ 处，边界条件现在是
\be\label{kilo}\psi_+=\Rf\psi_-,\ee
其中 $\Rf:TY\to TY$ 是反射，在 $L$ 的切丛 $TL$ 上作用为 $+1$，在其正交补 $TL_\perp$ 上作用为 $-1$。特别地 $\Rf^2=1$。合并两个条件，沿弦的费米零模必须满足 $J\psi_-=\Rf\psi_-$，或（因 $J^2=-1$）$J\Rf\psi_-=-\psi_-$。

我们先假设 $\omega_J$ 限制到 $L$ 上非退化，并证明此时 $-1$ 不是 $J\Rf$ 的特征值，从而沿弦没有费米零模。如第 \ref{abq} 节所解释，一旦知道这一点，玻色零模有效作用量的分析就化为 $L$ 的量子化问题。

为证明 $-1$ 不是 $J\Rf$ 的特征值，设 $J\Rf\alpha=-\alpha$，并写 $\alpha=\alpha_++\alpha_-$，其中 $\alpha_+\in TL$、$\alpha_-\in TL_\perp$。由 $J\Rf\alpha=-\alpha$ 与 $J^2=-1$ 可得 $J\alpha_+=-\alpha_-$、$J\alpha_-=\alpha_+$。对任意 $\chi\in TL$，有
$\omega_J(\chi,\alpha_+)=g(\chi,J\alpha_+)=-g(\chi,\alpha_-)=0$，其中 $\omega_J(\ ,\ )$ 与 $g(\ ,\ )$ 是 $\omega_J$ 与 $g$ 相伴的反对称与对称配对，而 $g(\chi,\alpha_-)=0$ 是因为 $\chi$ 取值于 $TL$、$\alpha_-$ 取值于 $TL_\perp$。因此，若 $\omega_J$ 限制到 $TL$ 上非退化，则 $\alpha_+$ 必须为零，于是 $\alpha_-=-J\alpha_+$ 也为零，$\alpha=0$。故 $J\Rf$ 没有特征值 $-1$，沿弦没有费米零模。

现在讨论相反情形：$\omega_J$ 限制到 $TL$ 上为零。这意味着对任意 $\chi,\chi'\in TL$，有 $0=\omega_J(\chi,\chi')=g(\chi,J\chi')$。换言之，$J$ 把 $TL$ 映到 $TL_\perp$；又因 $J^2=-1$，它把 $TL_\perp$ 映到 $TL$，且与 $\Rf$ 反交换。由此可得 $(J\Rf)^2=1$，且 $J\Rf$ 的特征值 $-1$ 的特征空间是中间维的，其到 $TL$ 的正交投影是同构。费米零模的量子化因此给出 $L$ 的 spin 丛的一个拷贝；对复流形 $L$，这可以等同于 $K^{1/2}\otimes\oplus_{q=1}^{\dim_\C L}\Omega^{0,q}(L)$，其中 $\Omega^{0,q}(L)$ 是 $L$ 上 $(0,q)$-形式丛。

玻色零模的变分仍如式 \eqref{btrans} 所示，但现在玻色零模描述到 $L$ 而非到 $Y$ 的映射，且 $\rho$ 对应于 $L$ 上的 $(0,1)$-形式。因此在这个描述层次上，$Q$ 的上同调化为 $L$ 上 $\bar\partial$ 算符的上同调。这个 $\bar\partial$ 算符作用在 $\L^{-1}\otimes K^{1/2}$ 的截面上，其中因子 $\L^{-1}$ 来自 $\B$ 的 CP 丛（$\B_\cc$ 的 CP 丛平凡，没有类似因子），因子 $K^{1/2}$ 来自费米零模的量子化。于是 $Q$ 的上同调为 $\oplus_qH^q(L,\L^{-1}\otimes K^{1/2})$，或简写为 $H^*(L,\L^{-1}\otimes K^{1/2})$。如正文所解释，$\L$ 与 $K^{1/2}$ 一般不是普通复线丛，而是 $\spinc$ 结构，但张量积 $\L^{-1}\otimes K^{1/2}$ 总是良定义的复线丛。

回到 $\Hom(\B_\cc,\B_\cc)$ 的情形，不必包含 $K^{1/2}$ 因子的理由很简单：复辛流形 $Y$ 的典则丛总是平凡的（它被全纯辛形式 $\Omega$ 的顶层外幂平凡化）。故 $Y$ 有一个 spin 丛就是 $\oplus_{q=0}^{\dim_\C Y}\Omega^{0,q}(Y)$。若 $Y$ 不是单连通的，它可能有其他 spin 丛，但 $\oplus_{q=0}^{\dim_\C Y}\Omega^{0,q}(Y)$ 是应当使用的那个，因为恒等算符作为 $\Hom(\B_\cc,\B_\cc)$ 的元素存在；若用 $Y$ 上不同的 spin 结构量子化费米零模，就会失去这一点。

\subsection{修正}\label{corrections}

最后讨论下面这个重要问题。本附录中我们做的所有计算都基于 $\sigma$-模型微扰论的最低阶；换言之，它们对应于 $\hbar$ 或 $1/\Omega$ 展开的首项。是否有重要的修正？

在 $\Hom(\B,\B_\cc)$ 的例子中（其中 $\omega_J$ 限制到 $L$ 上非退化），上述论证的结论是证明没有费米零模。故在低能描述中 $Q$ 可近似为 0，问题化为用某个预量子线丛量子化 $L$。除了某些有利情形（特别是存在具有某些良好性质的极化时），我们并不声称对所得到的量子 Hilbert 空间 $\Hom(\B_\cc,\B)$ 有显式描述。因此我们并未作出任何可能受修正影响的断言。

在 $\Hom(\B,\B_\cc)$ 的情形中（$\B$ 的支撑是复拉格朗日子流形），以及在 $\Hom(\B_\cc,\B_\cc)$ 的情形中，一般可以想象对小 $\hbar$ 极限有修正。这个问题究竟意味着什么其实相当微妙。我们聚焦于 $\Hom(\B,\B_\cc)$，但 $\Hom(\B_\cc,\B_\cc)$ 的情形可类似讨论。

说 $\Hom(\B,\B_\cc)$ 相对于最低阶答案被“修正”，一种可能的含义是：高阶修正使超荷 $Q$ 的上同调小于我们作为 $Q$ 的最低阶近似所找到的 $\bar\partial$ 算符的上同调。例如，对复拉格朗日子流形 $L$，可能有 $u\in H^0(L,K_L^{1/2})$、$\eta\in H^1(L,K_L^{1/2})$，二者都被 $\bar\partial$ 消灭；但 $\sigma$-模型微扰论的高阶修正可能导致 $Qu=\veps\eta$（$\veps$ 为小参数），从而把 $u$ 与 $\eta$ 都从上同调中移除。这类修正只有在 $H^i(L,K_L^{1/2})$ 对两个相邻整数 $i,i+1$ 都非零时才可能。

保证在这个意义上没有修正的一个简单条件是 $H^i(L,\L^{-1}\otimes K_L^{1/2})=0$（$i>0$）。本文考虑的许多例子确实如此。在任何这类情形中，$\Hom(\B,\B_\cc)$ 在某种意义上是未受修正的。

然而，当能在更精确的意义上说 $\Hom(\B,\B_\cc)$ 未受修正时，事情有趣得多。在解释这一点之前，先指出两类重要的例子：在其中可以假设 $\L^{-1}\otimes K_L^{1/2}$ 平凡，从而最低阶答案化为 $H^i(L,\O)$。进一步，我们特殊化到 $H^i(L,\O)=0$（$i>0$）的情形，于是最低阶答案实际上化为 $H^0(L,\O)$，即 $L$ 上全纯函数环。

第一类满足这些判据的例子：设 $L\cong\C^n$ 是某个复辛流形 $Y$ 的全纯极化的纤维。则 $K_L$ 平凡；又因 $L$ 单连通，平坦线丛 $\L$ 与平方根 $K_L^{1/2}$ 都平凡。最后，$H^i(\C^n,\O)=0$（$i>0$）也成立。

第二类例子：设 $L\subset Y\times Y'$ 是与 $Y$ 的全纯辛同胚相伴的对应。则 $L$ 作为复流形同构于 $Y$，故 $K_L$ 平凡，且可选 $\L$ 与 $K_L^{1/2}$ 也平凡。最后，若 $H^i(Y,\O)=0$（许多有趣的复辛流形如此，当然并非全部），则 $L$ 也有同样性质。

在这两类情形中，$\Hom(\B,\B_\cc)$ 在首阶近似下都化为 $L$ 上全纯函数空间。这个空间含有一个特殊的子空间，即不在无穷远增长的函数——常数函数。在正文对这些例子的讨论中，与常数函数相伴的角起了重要作用。所以让我们仔细讨论：在目前所作的假设之外，要确保存在与常数函数 1 相伴的 $\Hom(\B,\B_\cc)$ 的典则元素，还涉及什么。

考虑 $\sigma$-模型微扰论中会发生什么；我们假设它由小参数 $\veps$ 控制（若把 $Y$ 的度量按常数 $t$ 缩放，则可取 $\veps=1/t$）。在 $\veps$ 的最低阶，$Q$ 化为 $\bar\partial$，它消灭常数函数。假设在 $\veps^k$ 阶（某个 $k>0$）遇到第一个非零贡献，其形式为
\be\label{firstcor}Q\cdot1=\alpha,\ee
其中 $\alpha$ 是 $L$ 上的 $(0,1)$-形式，阶为 $\veps^k$。由于 $Q^2=0$，将有 $Q\alpha=0$；在 $\veps^k$ 阶，这条件恰好化为 $\bar\partial\alpha=0$。为补偿这一效应，选取 $w$ 使
\be\label{newcor}\bar\partial w=-\alpha.\ee
这样的 $w$ 存在，因为 $\alpha$ 是 $\bar\partial$-闭的且 $H^1(L,\O)=0$。选定这样的 $w$ 后，把 1 换成 $\psi_1=1+w$，就看到 $Q\psi_1$ 到 $\veps^{k+1}$ 阶为零。在更高阶，可能发现 $Q\psi_1$ 在 $\veps^r$ 阶（某个 $r>k$）非零。若如此，重复同样程序，对 $\psi_1$ 作进一步修正。

在这个程序的每一阶段，都要求解形如 $\bar\partial w=-\alpha$ 的方程。这样的 $w$ 总存在（因为 $\alpha$ 总是 $\bar\partial$-闭的），且 $w$ 在可加一个全纯函数的意义下唯一。若总能以典则方式选取 $w$，则该程序给出与常数函数 1 相伴的特殊 $\psi_1$。因此，并不是常数函数 1 “是” $\Hom(\B,\B_\cc)$ 的元素，而是存在一个与常数函数 1 典则相伴的 $\Hom(\B,\B_\cc)$ 元素。这个陈述正是应用中实际需要知道的。

如何指定方程 $\bar\partial w=-\alpha$ 的典则解？显然的办法是要求 $w$ 沿 $L$ 在无穷远消失。若方程有在无穷远消失的解，则该解唯一，因为任何其他解都可通过加一个全纯函数得到。若 $\alpha$ 在无穷远消失得足够快，则方程确有在无穷远消失的解。例如，对用 $z_1,\cdots,z_n$ 参数化、度量为 $\sum_{k=1}^n|\d z_k|^2$ 的 $L=\C^n$，若 $\alpha$ 在无穷远比 $1/|z|$ 消失得更快，则 $w$ 的方程有在无穷远消失的解。

作为一个良好情形的重要例子，设 $Y$ 与 $L$ 渐近锥形。这特别意味着它们的曲率在无穷远消失。由于对首阶答案的 $\sigma$-模型修正涉及 $L$ 的内蕴与外蕴曲率以及 $Y$ 的曲率，在这种情形中 $\alpha$ 会在无穷远消失得足够快，且可假设 $w$ 在无穷远消失。上述分析的结论有望作为精确陈述成立，尽管论证是在微扰论中表述的。

最后用几个简单例子说明潜在的陷阱。第一个例子：令 $Y=\C^2$，用 $p,q$ 参数化，辛形式为 $\d p\d q$，平坦超 Kähler 度量为 $|\d p|^2+|\d q|^2$。设 $Y'$ 是 $Y$ 的第二个拷贝，用 $p',q'$ 参数化，考虑由
\begin{align}\label{cordef}p'&=p+q^n,\cr q'&=q\end{align}
定义的拉格朗日对应 $L$。$Y\times Y'$ 的度量为 $|\d p|^2+|\d q|^2+|\d p'|^2+|\d q'|^2$。令 $q'=q$ 后，$L$ 的度量为 $|\d p|^2+|\d p'|^2+2|\d q|^2$，其中 $p,p',q$ 由 $p'=p+q^n$ 相联系。令 $p_\pm=(p'\pm p)/\sqrt2$。则 $L=\C\times L'$，其中 $\C$ 由 $p_+$ 参数化、度量为平坦的 $|\d p_+|^2$，而 $L'$ 由 $p_-,q$ 参数化，二者由
\be\label{fleq}q^n=\sqrt2 p_-\ee
相联系，度量为 $|\d p_-|^2+2|\d q|^2$。由于 $Y\times Y'$ 与 $L$ 有把 $p_+$ 平移常数的精确对称，$L=\C\times L'$ 的第一个因子从 $\sigma$-模型微扰论中解耦，上述讨论中的方程 $\bar\partial w=-\alpha$ 可理解为 $L'$ 上的方程。$L'$ 只有一个无穷“端”，即 $p_-,q\to\infty$。这个区域便于用 $p_-$ 参数化，此时 $q\sim p_-^{1/n}$；式 \eqref{fleq} 与 $L'$ 度量的形式表明 $L'$ 渐近平坦，曲率按 $1/|p_-|^{4-2/n}$ 消失。因此总有在无穷远消失的解 $w$；这是应当存在与常数函数 1 典则相伴的 $\Hom(\B,\B_\cc)$ 元素的例子。

$\C^{2n}$ 的驯顺辛同胚群的生成元（用第 \ref{oversym} 节的语言）对任意 $n$ 可类似处理。不过，让我们考虑 $\C^{2n}$ 的一个假想的非驯顺辛同胚。用坐标 $x^1,x^2,\cdots,x^{2n}$ 参数化 $\C^{2n}$，与 $\C^{2n}$ 的辛同胚相伴的对应 $L$ 一般由方程
\begin{align}\label{expfno}x'{}^i&=f^i(x^k),\cr x^j&=k^j(x'{}^k)\end{align}
描述。对一般的非驯顺辛同胚，关于函数 $f^i$ 与逆函数 $k^j$ 所知甚少，不清楚如何描述 $L$ 的无穷远区域，也不清楚如何论证方程 $\bar\partial w=-\alpha$ 有在无穷远消失的解。对这一点若有更多洞见，本可在第 \ref{symmeq} 节关于 $\C^{2n}$ 辛同胚的讨论中得到更强的结论。
